% Created 2026-03-16 Mon 07:59
% Intended LaTeX compiler: pdflatex
\documentclass[11pt]{article}
\usepackage[utf8]{inputenc}
\usepackage[T1]{fontenc}
\usepackage{graphicx}
\usepackage{longtable}
\usepackage{wrapfig}
\usepackage{rotating}
\usepackage[normalem]{ulem}
\usepackage{amsmath}
\usepackage{amssymb}
\usepackage{capt-of}
\usepackage{hyperref}
\author{Prologos Research}
\date{2026-03-16}
\title{Next-Generation Logic Programming: Beyond Prolog in a Propagator-First Architecture}
\hypersetup{
 pdfauthor={Prologos Research},
 pdftitle={Next-Generation Logic Programming: Beyond Prolog in a Propagator-First Architecture},
 pdfkeywords={},
 pdfsubject={},
 pdfcreator={Emacs 30.2 (Org mode 9.7.39)}, 
 pdflang={English}}
\begin{document}

\maketitle
\setcounter{tocdepth}{3}
\tableofcontents

\section{Introduction}
\label{sec:org8580cbc}

Prologos already has a relational sublanguage (\texttt{defr}, \texttt{rel}, \texttt{solve}, \texttt{explain})
backed by a persistent propagator network, bilattice-based well-founded semantics,
SLG-style tabling, ATMS-based assumption management, and narrowing. This document
explores how to extend this foundation beyond a Prolog port into a genuinely
next-generation logic programming system.

The key insight is that Prologos's propagator-first architecture is not merely an
implementation detail---it is a \emph{design space multiplier}. Where classical LP
systems must bolt on constraint solvers, probabilistic inference, or coinductive
reasoning as special-case extensions, a propagator network with pluggable lattices
can express all of these as instances of the same fixpoint computation.

This document surveys seven major extension axes and one integrative frontier:

\begin{enumerate}
\item Mode prefixes on logic variables (\texttt{+}, \texttt{-}, \texttt{?})
\item Domain-constrained logic variables (\texttt{?x:DomainType})
\item Co-inductive logic programming and greatest fixpoints
\item Rule learning and inductive logic programming
\item Probabilistic logic programming
\item Category-theoretic foundations
\item Novel directions (linear LP, temporal LP, epistemic LP, abduction, ASP, CHR, tabling)
\item The frontier: what a fresh redesign of LP looks like
\end{enumerate}

Each section connects back to Prologos's existing infrastructure with concrete
implementation considerations.
\subsection{Existing Infrastructure Reference}
\label{sec:org89bb739}

\begin{center}
\begin{tabular}{lll}
Module & Role & Key Abstractions\\
\hline
\texttt{propagator.rkt} & Persistent propagator network & \texttt{prop-cell}, \texttt{propagator}, \texttt{prop-network}, \texttt{cell-id}\\
\texttt{bilattice.rkt} & Three-valued reasoning (WFS) & \texttt{lattice-desc}, \texttt{bilattice-var}, \texttt{bool-lattice}\\
\texttt{atms.rkt} & Assumption-based truth maintenance & \texttt{assumption}, \texttt{supported-value}, \texttt{tms-cell}, \texttt{atms}\\
\texttt{tabling.rkt} & SLG-style memoization & \texttt{table-entry}, \texttt{table-store}, \texttt{wf-table-entry}\\
\texttt{wf-engine.rkt} & Well-founded semantics orchestration & \texttt{wf-answer}, \texttt{wf-naf-oracle}, bilattice fixpoint loop\\
\texttt{relations.rkt} & Relation registry and DFS solver & \texttt{relation-info}, \texttt{solve-goal}, \texttt{explain-goal}\\
\texttt{solver.rkt} & Solver configuration & \texttt{solver-config}, strategy/threshold/tabling/provenance\\
\texttt{narrowing.rkt} & Narrowing search (functional-logic bridge) & Value ordering, constraint propagation\\
\texttt{provenance.rkt} & Derivation chain tracking & Explanation trees, derivation depth\\
\end{tabular}
\end{center}
\section{Mode Prefixes on Logic Variables}
\label{sec:org58cbd66}

\subsection{Background: What Modes Are}
\label{sec:org78647e8}

In logic programming, a \emph{mode} declaration specifies the intended data-flow
pattern of a predicate's arguments:

\begin{itemize}
\item \texttt{+} (input): the argument is ground (fully instantiated) on entry
\item \texttt{-} (output): the argument is free on entry and will be bound by the predicate
\item \texttt{?} (unknown): the argument may be either ground or free
\end{itemize}

Mercury pioneered mandatory mode declarations with determinism inference. Ciao
Prolog uses a rich assertion language where modes are optional but enable powerful
static analysis via CiaoPP. Both systems demonstrate that modes unlock significant
optimization opportunities.
\subsection{Mercury's Mode System}
\label{sec:org9ba4e2d}

Mercury's execution algorithm exploits modes in several ways:

\begin{itemize}
\item \textbf{Determinism categories}: From modes, Mercury infers whether a predicate is
\texttt{det} (exactly one solution), \texttt{semidet} (zero or one), \texttt{multi} (one or more),
or \texttt{nondet} (zero or more). Deterministic predicates compile to simple function
calls with no backtracking machinery.
\item \textbf{Clause selection}: For \texttt{det=/=semidet} predicates with input arguments, the
compiler generates switch statements or hash lookups instead of sequential
clause search.
\item \textbf{Argument reordering}: The compiler may reorder goals within a clause body to
ensure input arguments are ground before they are needed, potentially converting
\texttt{nondet} goals to \texttt{semidet} ones.
\item \textbf{Specialization}: Multiple modes for the same predicate generate separate
compiled versions (e.g., \texttt{append(+, +, -)} vs \texttt{append(-, -, +)}).
\end{itemize}
\subsection{Ciao's Assertion Language}
\label{sec:org2e63603}

Ciao's approach is more flexible: modes are expressed as assertions that serve
both as documentation and as targets for static/dynamic verification. CiaoPP
infers types, modes, sharing, determinacy, non-failure, and resource bounds.

The assertion language supports properties like:
\begin{itemize}
\item \texttt{:- pred append(+list, +list, -list).} --- mode declaration
\item \texttt{:- check comp append(X, Y, Z) : (ground(X), ground(Y)) => ground(Z).} --- conditional mode
\end{itemize}
\subsection{Proposed Prologos Syntax}
\label{sec:orgdcf5a0b}

\begin{verbatim}
;; Mode-annotated relation
defr append {A : Type}
  | [+xs : List A] [+ys : List A] [-zs : List A]
  ;; Ground inputs, free output -- deterministic append

;; Multi-mode relation
defr member {A : Type}
  | [+x : A] [+xs : List A]         ;; membership test (semidet)
  | [-x : A] [+xs : List A]         ;; enumeration (nondet)

;; Solver can exploit mode information
solve [member ?who '[:alice :bob :carol]]
;; With ?who free and the list ground, the solver enumerates (nondet mode)

solve [member :bob +known-list]
;; With both ground, the solver does a membership test (semidet mode)
\end{verbatim}
\subsection{Determinism Inference from Modes}
\label{sec:org0c5c5b2}

\begin{center}
\begin{tabular}{llll}
Input Modes & Pattern & Determinism & Optimization\\
\hline
All \texttt{+} & First-match & \texttt{det} & Compile to function call\\
All \texttt{+} but last & Generate & \texttt{semidet} & Single-solution search, cut on success\\
Mix of \texttt{+}, \texttt{-} & Search & \texttt{multi} & Indexed on \texttt{+} args, enumerate \texttt{-} args\\
All \texttt{-} & Enumerate & \texttt{nondet} & Full backtracking search\\
\end{tabular}
\end{center}
\subsection{Composition with the Propagator Network}
\label{sec:org28d94c6}

Modes interact with Prologos's propagator network in two ways:

\begin{enumerate}
\item \textbf{Cell initialization}: A \texttt{+} argument maps to a propagator cell initialized
with a ground value (\texttt{net-cell-write} at construction). A \texttt{-} argument maps
to a cell initialized at \texttt{bot}. A \texttt{?} argument maps to a cell that may or
may not have an initial value.

\item \textbf{Propagation direction}: Mode analysis determines which propagators fire
first. In a clause like \texttt{append(+X, +Y, -Z)}, the propagator that computes
\texttt{Z} from \texttt{X} and \texttt{Y} should fire eagerly (forward propagation). In
\texttt{append(-X, -Y, +Z)}, the reverse propagators decomposing \texttt{Z} should fire.

\item \textbf{Determinism as lattice metadata}: The determinism category of a goal
(\texttt{det}, \texttt{semidet}, \texttt{multi}, \texttt{nondet}) can be tracked as metadata on the
propagator cell, allowing the scheduler (\texttt{run-to-quiescence}) to prioritize
deterministic propagations.
\end{enumerate}
\subsection{Implementation Roadmap}
\label{sec:org1e92642}

\begin{enumerate}
\item Parse mode prefixes (\texttt{+}, \texttt{-}, \texttt{?}) in \texttt{defr} parameter lists
\item Extend \texttt{param-info} in \texttt{relations.rkt} with a \texttt{mode} field
\item Add determinism inference in a new \texttt{mode-analysis.rkt} module
\item Modify \texttt{solve-goal} to select clause ordering and indexing based on modes
\item Generate specialized propagator networks per mode combination
\end{enumerate}
\section{Domain-Constrained Logic Variables}
\label{sec:orgb85bb10}

\subsection{Background: CLP(X) Family}
\label{sec:org5af9141}

Classical Prolog treats logic variables as untyped, unifiable terms. The CLP(X)
family parameterizes the solver over a constraint domain X:

\begin{center}
\begin{tabular}{llll}
Domain & Variables Over & Constraints & Solving Technique\\
\hline
CLP(FD) & Finite integers & \texttt{X \#> Y}, \texttt{X in 1..10} & Arc consistency, propagation\\
CLP(R) & Real numbers & \texttt{X > 3.14}, \texttt{X + Y =:} 5.0= & Simplex, interval arithmetic\\
CLP(Q) & Rationals & Same as CLP(R) & Exact rational arithmetic\\
CLP(Set) & Finite sets & \texttt{X subset Y}, \texttt{card(X) = 3} & Set-interval propagation\\
CLP(Bool) & Booleans & \texttt{X \#\textbackslash{}/ Y}, \texttt{X \#/\textbackslash{} Z} & BDD, SAT\\
CLP(S) & Strings & \texttt{concat(X, Y, Z)} & String constraint solvers\\
\end{tabular}
\end{center}
\subsection{Proposed Prologos Syntax}
\label{sec:org2147629}

\begin{verbatim}
;; Domain-constrained logic variable
defr schedule
  | [?room : FD 1..10] [?time : FD 0..23] [?course : Symbol]

;; Constraints use trait-based dispatch
rel schedule
  | [?r ?t :math101]  :- [?t > 8N] [?t < 17N]
  | [?r ?t :phys201]  :- [?t > 9N] [?t < 18N] [?r /= 5N]

;; Real-valued constraints
defr trajectory
  | [?x : Real] [?y : Real] [?t : Real 0.0..100.0]
  :- [?y = [?x * ?x]] [?t >= 0.0]
\end{verbatim}
\subsection{Arc Consistency and Bounds Propagation as Propagator Cells}
\label{sec:org92cc26b}

The key insight is that CLP constraint propagation \emph{is} propagator network
computation:

\begin{itemize}
\item An FD variable \texttt{?x : FD 1..10} maps to a propagator cell whose value is a
domain representation (e.g., a sorted interval list or bitset). The lattice
is the \emph{reverse inclusion order} on domains: removing values from the domain
is the monotone operation (domain shrinks over time).

\item Arc consistency for a binary constraint \texttt{c(X, Y)} is a propagator watching
\texttt{X} and \texttt{Y}: when \texttt{X}'s domain shrinks, it filters \texttt{Y}'s domain (and vice
versa). This is exactly a two-input, two-output propagator in the network.

\item Bounds propagation is a coarser version: the cell stores only \texttt{[lo, hi]} and
the propagator updates bounds. This maps to a simpler lattice (pairs of
integers with component-wise narrowing).
\end{itemize}
\subsection{Lattice Instances for Constraint Domains}
\label{sec:orgf6ec7e3}

These connect directly to the lattice catalog (\texttt{docs/research/2026-03-14\_LATTICE\_CATALOG.org}):

\begin{center}
\begin{tabular}{lllll}
Domain & Cell Lattice & Merge (\texttt{join}) & Bot & Top\\
\hline
FD & Reverse powerset & Intersection & Full domain & Empty set (=!)\\
Bounds & Interval \texttt{\textbackslash{}[lo, hi\textbackslash{}]} & Narrow \texttt{[max lo, min hi]} & \texttt{[-inf, +inf]} & Empty interval\\
Real & Interval arithmetic & Narrow & \texttt{(-inf, +inf)} & Empty\\
Set & Set intervals \texttt{[lb, ub]} & Narrow both bounds & \texttt{[empty, U]} & \texttt{lb > ub}\\
Boolean & Trilean \texttt{\{T, F, unknown\}} & Refine & \texttt{unknown} & Contradiction\\
\end{tabular}
\end{center}
\subsection{Composition with ATMS}
\label{sec:org76d9b68}

Domain constraints compose with the ATMS (\texttt{atms.rkt}) for speculative reasoning:

\begin{verbatim}
;; Under different assumptions, a variable may have different domains
solve [schedule ?room ?time :math101]
  assuming { room-available 3 }
  ;; ATMS tracks: {room-available 3} => room domain shrinks to {3}
  ;; Other worldviews may have different domain states
\end{verbatim}

Each \texttt{supported-value} in a \texttt{tms-cell} can carry a domain state. The ATMS
manages which domain states are consistent with which assumption sets,
enabling the solver to explore multiple constraint scenarios simultaneously.
\subsection{Implementation Roadmap}
\label{sec:org3827cf6}

\begin{enumerate}
\item Define \texttt{domain-cell} constructors in \texttt{propagator.rkt} (FD, bounds, interval)
\item Implement constraint propagators (arc-consistency, bounds-propagation)
\item Extend \texttt{?x:Type} syntax in the relational parser to accept domain annotations
\item Add domain-aware unification in the solver (\texttt{unify} + domain narrowing)
\item Connect to the existing \texttt{narrowing.rkt} infrastructure
\end{enumerate}
\section{Co-Inductive Logic Programming and Greatest Fixpoints}
\label{sec:orgaa8daef}

\subsection{Background: Coinduction as Dual of Induction}
\label{sec:org287422e}

Standard logic programming computes \emph{least fixpoints} (lfp): start from nothing,
build up derivable facts by applying clauses bottom-up until no new facts emerge.
Co-inductive logic programming computes \emph{greatest fixpoints} (gfp): start from
everything, eliminate non-derivable facts top-down.

The duality:
\begin{itemize}
\item \textbf{Induction} (lfp): finite derivation trees, well-founded recursion, Herbrand models
\item \textbf{Coinduction} (gfp): potentially infinite proof trees, corecursion, regular trees
\end{itemize}
\subsection{CoLP: Coinductive Hypothesis Rule}
\label{sec:orgb20be60}

Gupta, Bansal, and Simon (2007) introduced CoLP with the \emph{coinductive hypothesis
rule}: when proving a goal \texttt{G}, if \texttt{G} already appears as an ancestor in the
current proof search, succeed (coinductive hypothesis) rather than failing
(inductive loop detection).

This enables:

\begin{verbatim}
;; Coinductive definition of infinite streams
:coinductive
defr stream-of
  | [?s : Stream A] [?a : A]

rel stream-of
  | [?s ?a] :- [head ?s ?a] [tail ?s ?rest] [stream-of ?rest ?a]
  ;; Under coinduction, stream-of [ones 1N] succeeds for ones = 1 :: ones

;; Bisimulation as coinductive equality
:coinductive
defr bisimilar
  | [?s1 : Stream A] [?s2 : Stream A]

rel bisimilar
  | [?s1 ?s2] :- [head ?s1 ?h] [head ?s2 ?h]
                  [tail ?s1 ?t1] [tail ?s2 ?t2]
                  [bisimilar ?t1 ?t2]
\end{verbatim}
\subsection{Flexible Coinduction: Coclauses}
\label{sec:org1366910}

Ancona, Barbieri, and Zucca (2020) introduced \emph{coclauses} that allow tuning the
coinductive behavior per-predicate. A coclause is matched when a coinductive
hypothesis fires, allowing the user to specify \emph{what} should happen when a loop
is detected rather than always succeeding.

\begin{verbatim}
;; Flexible coinduction with coclause
:coinductive
defr subtype
  | [?a : Type] [?b : Type]

;; Regular clauses (inductive base cases)
rel subtype
  | [:Nat :Int]
  | [[List ?a] [List ?b]] :- [subtype ?a ?b]

;; Coclause: when coinductive hypothesis fires, check structural compatibility
coclause subtype
  | [?a ?b] :- [structural-compat ?a ?b]
\end{verbatim}
\subsection{Composition with Bilattice (Well-Founded Engine)}
\label{sec:org01b6ede}

The existing bilattice infrastructure in \texttt{bilattice.rkt} already provides the
dual ascending/descending cell pairs needed for coinductive reasoning:

\begin{itemize}
\item \textbf{Ascending cell} (lfp direction): accumulates evidence FOR a predicate being
true. This is standard inductive LP.
\item \textbf{Descending cell} (gfp direction): starts at \texttt{top} and eliminates evidence
AGAINST a predicate being true. This is the coinductive direction.
\end{itemize}

The bilattice reading at quiescence gives exactly the three-valued approximation
needed:
\begin{itemize}
\item \texttt{lower = upper = true}: predicate is both inductively and coinductively true
\item \texttt{lower = false, upper = true}: unknown (gap) --- the predicate might be true
coinductively but has no finite inductive proof
\item \texttt{lower = upper = false}: definitively false even coinductively
\end{itemize}

The \texttt{wf-engine.rkt} iterative fixpoint loop already alternates between bottom-up
and top-down phases. Extending it with coinductive hypothesis detection requires:

\begin{enumerate}
\item Maintaining a \emph{coinductive stack} during DFS proof search
\item When a goal matches an ancestor, consult the bilattice upper cell
\item If upper cell is \texttt{true} (not yet refuted), succeed coinductively
\item The refutation phase may later push the upper cell to \texttt{false}
\end{enumerate}
\subsection{Coinductive Session Types}
\label{sec:org11fad71}

Prologos has session types. Session types are naturally coinductive (a server
loop is an infinite protocol). Coinductive LP enables:

\begin{verbatim}
;; Session type as coinductive predicate
:coinductive
defr conforms
  | [?process : Process] [?session : Session]

rel conforms
  | [?p ?s] :- [step ?p ?action ?p'] [transition ?s ?action ?s']
                [conforms ?p' ?s']
;; Coinductive: the process conforms if it can always make a matching step
\end{verbatim}
\subsection{Implementation Roadmap}
\label{sec:orgf5052c0}

\begin{enumerate}
\item Add \texttt{:coinductive} annotation to \texttt{defr} (in \texttt{relation-info}, new field)
\item Track coinductive ancestor stack in the DFS solver
\item Connect to bilattice upper cell for coinductive hypothesis resolution
\item Implement coclause syntax and matching
\item Test with session type conformance checking
\end{enumerate}
\section{Rule Learning and Inductive Logic Programming}
\label{sec:orgf2296af}

\subsection{Background: ILP Systems}
\label{sec:org0b9bc2e}

Inductive Logic Programming (ILP) learns logic programs (rules) from examples
and background knowledge. The major modern systems:

\begin{center}
\begin{tabular}{llll}
System & Method & Language & Key Feature\\
\hline
Metagol & Meta-interpretive & Prolog & Metarule-guided search, predicate invention\\
ILASP & ASP-based & ASP & Optimal w.r.t. cost functions\\
Popper & Hypothesis testing & Prolog + ASP & Generate-test-constrain loop\\
Louise & Polynomial MIL & Prolog & Scalable meta-interpretive learning\\
FastLAS & Fast ILASP variant & ASP & Noise-tolerant\\
\end{tabular}
\end{center}
\subsection{Meta-Interpretive Learning (MIL)}
\label{sec:org94f630a}

MIL learns by specializing \emph{metarules}---higher-order clause templates:

\begin{verbatim}
%% Metarule: chain
metarule(chain, [P,Q,R], [P,A,B] :- [[Q,A,C],[R,C,B]]).
%% "P(A,B) :- Q(A,C), R(C,B)" — learn transitive compositions
\end{verbatim}

The hypothesis space is controlled by the set of available metarules.
Typed MIL (MetagolT) uses type checking to prune the search space by a
cubic factor.
\subsection{Proposed Prologos Surface Syntax for ILP}
\label{sec:orgf0fdc76}

\begin{verbatim}
;; Define a learning task
learn grandparent
  from-examples
    [grandparent :alice :carol] ;; positive
    [grandparent :alice :dave]  ;; positive
  negative-examples
    [grandparent :alice :alice] ;; negative
  background
    parent member
  metarules
    chain identity
  max-clauses 3

;; The system would synthesize:
;; rel grandparent | [?x ?z] :- [parent ?x ?y] [parent ?y ?z]
\end{verbatim}
\subsection{Differentiable ILP}
\label{sec:org295b23f}

Recent work (NeuralLP, DRUM, and GLIDR) embeds ILP in differentiable
frameworks. Key ideas:

\begin{itemize}
\item Clause selection is parameterized by continuous weights (softmax over metarules)
\item Forward chaining is differentiable (matrix operations over relation adjacency)
\item Gradient descent optimizes rule weights against example loss
\end{itemize}

This connects to Prologos through the propagator network:

\begin{itemize}
\item Each candidate clause becomes a propagator with a continuous weight
\item The weight lives in a propagator cell over a \texttt{\textbackslash{}[0,1\textbackslash{}]} lattice
\item Gradient information flows backward through the network (counter-propagation)
\item The ATMS tracks which clause combinations are consistent
\end{itemize}

\begin{verbatim}
;; Differentiable ILP as weighted clauses
learn-differentiable path
  from-data edge-triples
  metarules [chain identity inverse]
  optimizer :adam
  epochs 100
  ;; Each learned clause gets a weight:
  ;; rel path | [?x ?y] :- [edge ?x ?y]        ;; weight 0.95
  ;; rel path | [?x ?z] :- [edge ?x ?y] [path ?y ?z]  ;; weight 0.88
\end{verbatim}
\subsection{Composition with Propagator Networks}
\label{sec:orgfafff6b}

ILP as propagator computation:

\begin{enumerate}
\item \textbf{Hypothesis space as lattice}: The set of candidate programs, ordered by
generality (more general = higher). Learning descends from the most general
hypothesis.
\item \textbf{Example propagators}: Each positive/negative example is a propagator that
constrains the hypothesis space (positive: must cover; negative: must not cover).
\item \textbf{Metarule propagators}: Each metarule instantiation is a propagator that
adds a candidate clause to the hypothesis.
\item \textbf{Pruning propagators}: Constraint propagation eliminates inconsistent
hypothesis subspaces.
\end{enumerate}

The ATMS is particularly valuable here: each candidate clause is an assumption,
and nogoods track clause combinations that violate negative examples.
\subsection{Implementation Roadmap}
\label{sec:orgdf3a6c4}

\begin{enumerate}
\item Define \texttt{learn} top-level form and metarule syntax
\item Implement metarule-guided search as propagator network construction
\item Connect to ATMS for hypothesis tracking (clauses as assumptions)
\item Add differentiable ILP as a solver strategy (\texttt{:optimizer} config)
\item Bridge to the type system (typed MIL uses Prologos's type inference)
\end{enumerate}
\section{Probabilistic Logic Programming}
\label{sec:org77a6f8c}

\subsection{Background: ProbLog and the Distribution Semantics}
\label{sec:org3ef8905}

ProbLog (KU Leuven) extends Prolog with probabilistic facts:

\begin{verbatim}
0.3 :: burglary.
0.1 :: earthquake.
alarm :- burglary.
alarm :- earthquake.
\end{verbatim}

The \emph{distribution semantics}: each probabilistic fact independently true/false
with its stated probability. A ground query's probability is the sum over all
possible worlds where it is derivable, weighted by world probability.

Inference uses \emph{knowledge compilation}: the derivation structure is compiled to
a Boolean formula (SDD, BDD, or d-DNNF), then weighted model counting computes
the probability.
\subsection{Algebraic ProbLog and Semirings}
\label{sec:orgc61edfb}

Algebraic ProbLog (aProbLog) generalizes from probabilities to arbitrary
commutative semirings. A semiring \texttt{(S, +, x, 0, 1)} replaces:

\begin{center}
\begin{tabular}{llll}
Semiring & \texttt{+} (combine) & \texttt{x} (chain) & Meaning\\
\hline
Probability & \texttt{+} & \texttt{x} & Standard probabilistic LP\\
Viterbi & \texttt{max} & \texttt{x} & Most probable proof\\
Tropical & \texttt{min} & \texttt{+} & Shortest path\\
Boolean & \texttt{or} & \texttt{and} & Standard LP\\
Fuzzy & \texttt{max} & \texttt{min} & Fuzzy LP\\
Gradient & \texttt{+} & \texttt{x} + chain rule & Differentiable inference\\
Access control & \texttt{meet} & \texttt{join} & Security lattice LP\\
\end{tabular}
\end{center}

This is directly relevant to Prologos: \emph{every semiring-based LP is a
propagator network with a specific lattice}.
\subsection{DeepProbLog: Neural Predicates}
\label{sec:orgfaee791}

DeepProbLog (Manhaeve et al., 2018, with continued development through 2025)
extends ProbLog with \emph{neural predicates}: predicates whose truth probability
is parameterized by a neural network.

\begin{verbatim}
nn(digit_net, [X], Y, [0..9]) :: digit(X, Y).
addition(X, Y, Z) :- digit(X, A), digit(Y, B), Z is A + B.
\end{verbatim}

The neural network's output is a probability distribution over the predicate's
possible values. Learning jointly optimizes the neural network weights and the
logical program structure.
\subsection{DC-ProbLog: Continuous Distributions (2024)}
\label{sec:org5255144}

DC-ProbLog (Zuidberghe et al., 2024) extends ProbLog to hybrid discrete-continuous
domains using \emph{distributional clauses}. The key contribution is a \emph{measure
semantics} that generalizes the distribution semantics to continuous random
variables, along with an inference engine (IALW---infinitesimal algebraic
likelihood weighting) based on knowledge compilation.
\subsection{Proposed Prologos Syntax}
\label{sec:orgc82a71b}

\begin{verbatim}
;; Probabilistic facts (discrete)
0.3 :: burglary
0.1 :: earthquake

;; Weighted logic variable
defr diagnosis
  | [?symptom : Symbol] [?disease : Symbol ^?p]
  ;; ?p is the marginal probability of ?disease given ?symptom

;; Distributional clause (continuous)
defr sensor-reading
  | [?sensor : Symbol] [?value : Real ~ Normal ?mu ?sigma]

rel sensor-reading
  | [:thermometer ?v] :- [?v ~ Normal 20.0 2.0]
  | [:barometer ?v]   :- [?v ~ Normal 1013.0 5.0]

;; Query with probability
solve [diagnosis :fever ?d] :probabilistic
;; Returns: {d: :flu, p: 0.7}, {d: :cold, p: 0.25}, ...

;; Conditional probability
solve [alarm] given [earthquake] :probabilistic
\end{verbatim}
\subsection{Semiring Framework as Propagator Lattice}
\label{sec:org382697a}

The semiring-based framework maps directly to the propagator network:

\begin{itemize}
\item A \emph{probability cell} is a propagator cell over the semiring \texttt{(R+, +, x, 0, 1)}.
The merge operation is semiring addition (combining evidence from multiple
derivation paths).

\item A \emph{Viterbi cell} uses \texttt{(R+, max, x, 0, 1)} for most-probable-explanation
queries. Same propagator structure, different lattice.

\item The \emph{gradient semiring} enables backpropagation through the logic program.
Each cell stores both a value and its gradient. The merge operation combines
values additively and gradients via the chain rule.
\end{itemize}

New propagator cell types needed:

\begin{center}
\begin{tabular}{llll}
Cell Type & Lattice & Merge & Use Case\\
\hline
\texttt{prob-cell} & \texttt{[0,1]} with \texttt{+} & Sum of proofs & Marginal probability\\
\texttt{viterbi-cell} & \texttt{[0,1]} with \texttt{max} & Best proof & MAP inference\\
\texttt{log-prob-cell} & \texttt{(-inf,0]} with \texttt{logadd} & Numerically stable & Large-scale inference\\
\texttt{gradient-cell} & \texttt{R x R} (value, gradient) & Sum + chain rule & Differentiable LP\\
\texttt{interval-cell} & Interval \texttt{[lo,hi]} & Narrow & Bounded probability\\
\texttt{distribution-cell} & Measure space & Mixture & Continuous distributions\\
\end{tabular}
\end{center}
\subsection{Knowledge Compilation and the ATMS}
\label{sec:orgc93aa8e}

ProbLog's inference via knowledge compilation (SDDs, d-DNNF) has a natural
counterpart in the ATMS (\texttt{atms.rkt}):

\begin{itemize}
\item Each probabilistic fact is an \emph{assumption} in the ATMS
\item Each derivation path is a \emph{support set} (conjunction of assumptions)
\item The probability of a query is the weighted count over consistent support sets
\item \texttt{atms-solve-all} already enumerates all consistent worldviews
\item The probability computation is a post-processing step: sum over worldviews,
weighted by product of assumption probabilities
\end{itemize}

This means probabilistic LP can be implemented as a thin layer over the
existing ATMS, without a separate knowledge compilation pipeline.
\subsection{Implementation Roadmap}
\label{sec:org82258a4}

\begin{enumerate}
\item Add probability annotations to facts (\texttt{0.3 :: fact}) in the parser
\item Implement \texttt{prob-cell} and \texttt{viterbi-cell} lattice types
\item Extend \texttt{atms-solve-all} with weighted enumeration
\item Add \texttt{:probabilistic} solver strategy in \texttt{solver-config}
\item Implement distributional clause syntax and continuous sampling
\item Bridge to the gradient semiring for DeepProbLog-style learning
\end{enumerate}
\section{Category-Theoretic Foundations}
\label{sec:org3764850}

\subsection{Why Category Theory Matters for Prologos}
\label{sec:orga42d8b6}

Prologos already uses categorical structures implicitly:

\begin{itemize}
\item Propagator cells form a \emph{category of lattices} with monotone maps as morphisms
\item The ATMS manages a \emph{fibered structure} where each worldview is a fiber
\item Well-founded semantics uses \emph{approximation fixpoint theory} (Denecker et al.),
which is fundamentally about operators on bilattices
\item Dependent types form a \emph{locally cartesian closed category}
\item Session types form a \emph{category of processes} with session type morphisms
\end{itemize}

Making these connections explicit enables principled extension and composition.
\subsection{Toposes as Generalized Logic}
\label{sec:org597b17b}

A topos is a category that behaves like a generalized universe of sets. Every
topos has an \emph{internal logic} that may be intuitionistic, classical, or
something more exotic. The subobject classifier \texttt{Omega} in a topos generalizes
the Boolean truth values \texttt{\{true, false\}} to an arbitrary Heyting algebra.

For Prologos:
\begin{itemize}
\item The standard Herbrand model is the topos \texttt{Set} (classical logic)
\item Three-valued logic (well-founded semantics) corresponds to the topos
\texttt{Set\textasciicircum{}\{3\}} or equivalently presheaves on the three-element chain
\item Probabilistic logic corresponds to the \emph{Giry monad} on \texttt{Meas} (the category
of measurable spaces)
\item Fuzzy logic corresponds to sheaves valued in \texttt{[0,1]}
\end{itemize}

A topos-aware logic engine could parameterize over the choice of topos, enabling
the same syntax to express classical, intuitionistic, fuzzy, or probabilistic
reasoning by swapping the underlying topos.

\begin{verbatim}
;; Topos-parameterized relation
defr reachable {T : Topos}
  | [?x : Node] [?y : Node] [:via T]

;; Classical: standard LP
solve [reachable :a :b :via Classical]

;; Three-valued: well-founded semantics (already implemented!)
solve [reachable :a :b :via WellFounded]

;; Fuzzy: degree-of-truth in [0,1]
solve [reachable :a :b :via Fuzzy]
\end{verbatim}
\subsection{Sheaves for Contextual Reasoning}
\label{sec:org130c981}

Sheaves generalize the idea of "data varying over a space" with \emph{gluing
conditions} that ensure local data can be consistently assembled into
global data. In programming language terms:

\begin{itemize}
\item A presheaf assigns data to each "context" (open set)
\item A sheaf adds the condition that compatible local data determines unique global data
\item The \emph{sheaf condition} is a coherence requirement
\end{itemize}

For Prologos, sheaves enable \emph{contextual logic programming}:

\begin{verbatim}
;; Facts that vary by context (location, time, agent)
defr temperature
  | [?loc : Location] [?temp : Real] [:context ?ctx]

;; Sheaf condition: overlapping contexts must agree
;; If ctx1 and ctx2 overlap, temperature readings must be consistent

;; Query with context restriction
solve [temperature :kitchen ?t :context :morning]
\end{verbatim}

The propagator network implements the sheaf gluing condition: consistency
propagators between overlapping contexts ensure the sheaf axiom holds.
\subsection{Kan Extensions for Program Transformation}
\label{sec:org741d267}

Kan extensions formalize the notion of "best approximation" when changing
the domain of a functor. In programming:

\begin{itemize}
\item \emph{Left Kan extension} corresponds to existential quantification / colimit /
"the best summary"
\item \emph{Right Kan extension} corresponds to universal quantification / limit /
"the best conservative approximation"
\end{itemize}

For Prologos, Kan extensions relate to:

\begin{enumerate}
\item \textbf{Abstract interpretation}: The abstraction function \texttt{alpha} is a left adjoint
(left Kan extension), and the concretization \texttt{gamma} is a right adjoint.
This is a Galois connection.

\item \textbf{Program optimization}: The continuation-passing-style transformation
corresponds to the codensity monad, which arises from a right Kan extension.
This connects to Prologos's narrowing infrastructure.

\item \textbf{Query optimization}: Magic set transformations (converting bottom-up
evaluation to goal-directed) can be expressed as Kan extensions along
the inclusion functor from relevant facts to all facts.
\end{enumerate}
\subsection{Galois Connections for Abstract Interpretation}
\label{sec:org058130a}

A Galois connection \texttt{<alpha, C, D, gamma>} between a concrete domain \texttt{C} and an
abstract domain \texttt{D} provides:

\begin{itemize}
\item \texttt{alpha: C -> D} (abstraction): maps concrete values to their best abstract approximation
\item \texttt{gamma: D -> C} (concretization): maps abstract values to the concrete values they represent
\item Soundness guarantee: \texttt{alpha . gamma <} id\textsubscript{D}= and \texttt{id\_C <} gamma . alpha=
\end{itemize}

In Prologos's propagator framework:

\begin{verbatim}
;; Abstract interpretation as cross-domain propagation
;; (already have net-add-cross-domain-propagator in propagator.rkt!)

;; Concrete domain: exact integer sets
;; Abstract domain: sign lattice {neg, zero, pos, top, bot}

;; The Galois connection is a pair of cross-domain propagators:
;; alpha: integer-set-cell -> sign-cell (abstraction propagator)
;; gamma: sign-cell -> integer-set-cell (concretization propagator)
\end{verbatim}

The existing \texttt{net-add-cross-domain-propagator} in \texttt{propagator.rkt} is already
the infrastructure needed for Galois-connected abstract interpretation.
\subsection{Fibered Categories for Dependent Types + Logic}
\label{sec:org91f88e2}

Prologos has both dependent types and logic programming. The categorical
framework that unifies these is \emph{fibered category theory} (Jacobs, 1999):

\begin{itemize}
\item A \emph{fibration} \texttt{p: E -> B} maps a "total" category \texttt{E} to a "base" category \texttt{B}
\item Each object \texttt{b} in \texttt{B} has a \emph{fiber} \texttt{E\_b} (the category of types/propositions
depending on \texttt{b})
\item Reindexing along morphisms in \texttt{B} gives substitution
\item The Curry-Howard-Lambek correspondence extends: types are objects in fibers,
terms are morphisms, dependent types are fibered structures
\end{itemize}

For Prologos, this means:
\begin{itemize}
\item The type \texttt{<(x : Nat) -> Vec x A>} lives in a fiber over \texttt{Nat}
\item A relational query \texttt{solve [vec-append ?xs ?ys ?zs]} where \texttt{?xs : Vec n A}
is a morphism in the fiber over \texttt{n}
\item Substitution of \texttt{n :} 3= reindexes the fiber, specializing the query
\end{itemize}
\subsection{Institutions for Multi-Logic Reasoning}
\label{sec:org92903c2}

Institutions (Goguen and Burstall, 1992) formalize the notion of a "logical
system" abstractly, enabling reasoning across multiple logics:

\begin{itemize}
\item An institution \texttt{I = (Sig, Sen, Mod, |})= specifies:
\begin{itemize}
\item \texttt{Sig}: category of signatures
\item \texttt{Sen}: functor from signatures to sentences
\item \texttt{Mod}: functor from signatures to models (contravariant)
\item \texttt{|=}: satisfaction relation (parameterized by signature)
\end{itemize}
\end{itemize}

For Prologos, institutions enable multi-logic queries:

\begin{verbatim}
;; A query that mixes classical and well-founded reasoning
solve
  [parent :alice ?x]                    ;; classical LP
  :and
  (wf [reachable ?x :goal])            ;; well-founded semantics
  :and
  (prob [risky ?x] > 0.5)              ;; probabilistic LP
\end{verbatim}

Each sub-query runs in its own "institution" (logic), and the institution
morphisms (signature translations) ensure consistency across the boundary.
The propagator network is the \emph{universal medium}: each institution maps to
a specific lattice/propagator configuration, and cross-domain propagators
implement the institution morphisms.
\subsection{Implementation Roadmap}
\label{sec:orga829ec4}

\begin{enumerate}
\item Make lattice descriptors (\texttt{lattice-desc}) a first-class abstraction
parameterizable by topos choice
\item Implement sheaf gluing as consistency propagators
\item Express abstract interpretation via cross-domain propagators + Galois connections
\item Document the fibered structure of dependent types + relational sublanguage
\item Design institution-aware multi-logic solver dispatch
\end{enumerate}
\section{Novel Directions}
\label{sec:orgfdbb8f3}

\subsection{Linear Logic Programming}
\label{sec:org49d0c8e}

Prologos already has QTT (Quantitative Type Theory) tracking resource usage
via multiplicities (\texttt{m0}, \texttt{m1}, \texttt{mw}). Extending this to the relational
sublanguage enables \emph{linear logic programming} where facts are \emph{consumed}
by derivation.

\begin{verbatim}
;; Linear fact: consumed when used
defr has-item {1}  ;; multiplicity 1 = linear
  | [?agent : Agent] [?item : Item]

;; Using a linear fact consumes it
rel craft
  | [?agent :sword] :- {1} [has-item ?agent :iron]
                        {1} [has-item ?agent :wood]
  ;; After this derivation, agent no longer has iron or wood

;; Unrestricted fact: can be used any number of times
defr knows {w}
  | [?agent : Agent] [?recipe : Recipe]
\end{verbatim}

The propagator network needs a new cell type for linear resources: a
\emph{counting cell} where the merge operation is addition (resource
accumulation) and consumption is subtraction (resource use). The lattice
is \texttt{Z} (integers) with a "non-negative" invariant.

Recent work (2024) on deadlock-free separation logic and the Linear Session
Abstract Machine (SAM) shows how linear logic, session types, and concurrency
compose, which is directly relevant to Prologos's session type system.
\subsection{Temporal Logic Programming}
\label{sec:org9853e94}

Express LTL (Linear Temporal Logic) and CTL (Computation Tree Logic)
properties as relational queries.

\begin{verbatim}
;; Temporal operators as relational combinators
defr eventually
  | [?prop : Prop] [?trace : Trace]

rel eventually
  | [?p ?t] :- [holds ?p [head ?t]]
  | [?p ?t] :- [eventually ?p [tail ?t]]

defr always
  | [?prop : Prop] [?trace : Trace]

:coinductive  ;; Greatest fixpoint!
rel always
  | [?p ?t] :- [holds ?p [head ?t]] [always ?p [tail ?t]]

;; CTL: exists-path and forall-path over branching traces
defr exists-eventually
  | [?prop : Prop] [?state : State]

rel exists-eventually
  | [?p ?s] :- [holds ?p ?s]
  | [?p ?s] :- [transition ?s ?s'] [exists-eventually ?p ?s']
\end{verbatim}

Note that \texttt{always} is naturally coinductive (gfp), while \texttt{eventually} is
inductive (lfp). The bilattice infrastructure handles both directions
simultaneously.

Temporal LP composes with the propagator network through \emph{trace cells}:
a cell whose lattice is a prefix ordering on execution traces. The
\texttt{run-to-quiescence} scheduler already handles iterative fixpoints over
such lattices.
\subsection{Epistemic Logic Programming}
\label{sec:orgf0e7172}

Knowledge and belief operators for multi-agent reasoning:

\begin{verbatim}
;; Epistemic operators
defr knows
  | [?agent : Agent] [?prop : Prop]

defr believes
  | [?agent : Agent] [?prop : Prop]

;; Common knowledge
defr common-knowledge
  | [?group : List Agent] [?prop : Prop]

;; Epistemic rules
rel knows
  | [:alice [safe ?code]] :- [observed :alice ?code]
  | [?a ?p]               :- [told ?a ?p ?source] [trusts ?a ?source]

;; Negative introspection: if alice doesn't know p, she knows she doesn't
rel knows
  | [?a [not-known ?p]] :- (not [knows ?a ?p])
\end{verbatim}

Each agent's knowledge state maps to a separate \emph{worldview} in the ATMS.
The ATMS assumption sets model possible worlds: agent \texttt{A} knows \texttt{P} iff
\texttt{P} holds in all worldviews consistent with \texttt{A}'s observations.
\subsection{Abductive Logic Programming}
\label{sec:org3a4a238}

Abduction explains observations by hypothesizing causes:

\begin{verbatim}
;; Abducible predicates (can be hypothesized)
:abducible
defr broken | [?component : Component]

:abducible
defr leak | [?pipe : Pipe]

;; Rules connecting causes to effects
rel alarm
  | [] :- [broken :sensor]
  | [] :- [leak :main-pipe] [not [valve-closed :main]]

;; Abductive query: what explains the observation?
abduce [alarm]
;; Returns: {broken :sensor} OR {leak :main-pipe, not valve-closed :main}
\end{verbatim}

The ATMS is the natural implementation vehicle: each abducible fact becomes
an assumption. \texttt{atms-solve-all} enumerates minimal consistent explanations.
The \texttt{atms-minimal-diagnoses} function already implements GDE-2 style
minimal diagnosis computation, which is precisely abductive reasoning.
\subsection{Answer Set Programming (ASP) Integration}
\label{sec:org4b03321}

ASP extends LP with stable model semantics, choice rules, and aggregates:

\begin{verbatim}
;; Choice rule: non-deterministically include/exclude
choice [color ?node ?c] :- [node ?node] [c in '[:red :green :blue]]

;; Integrity constraint: no two adjacent nodes same color
:- [edge ?n1 ?n2] [color ?n1 ?c] [color ?n2 ?c]

;; Aggregate
defr total-cost | [?cost : Nat]
rel total-cost
  | [?total] :- [?total = #sum { ?c : [task ?t] [cost ?t ?c] }]
\end{verbatim}

ASP's stable models correspond to the ATMS's consistent worldviews:
\begin{itemize}
\item Each choice rule creates an \texttt{atms-amb} (mutually exclusive assumptions)
\item Integrity constraints become \texttt{atms-add-nogood} (forbidden assumption sets)
\item \texttt{atms-solve-all} enumerates stable models
\end{itemize}

The connection between ASP and CHR has been established (Abdennadher et al.,
2021): first-order ASP programs can be compiled to CHR programs, eliminating
the grounding phase.
\subsection{Constraint Handling Rules (CHR)}
\label{sec:orgaccfc04}

CHR is a committed-choice rule language for implementing constraint solvers:

\begin{verbatim}
;; CHR rules for a simple constraint solver
;; Simplification: replace constraint with simpler one
chr leq-reflexive
  | [leq ?x ?x] <=> :true

;; Propagation: add new constraint without removing old
chr leq-transitive
  | [leq ?x ?y] [leq ?y ?z] ==> [leq ?x ?z]

;; Simpagation: some kept, some removed
chr leq-antisymmetric
  | [leq ?x ?y] \ [leq ?y ?x] <=> [?x = ?y]
\end{verbatim}

CHR maps directly to propagator networks:
\begin{itemize}
\item Each constraint is a propagator cell
\item Simplification rules are propagators that replace cell values
\item Propagation rules add new cells and propagators
\item The CHR store is the propagator network state
\end{itemize}

Recent work on Quantified CHR (QCHR, 2025) extends CHR with dynamic
quantifier generation, which could interact with Prologos's dependent
types.
\subsection{Advanced Tabling: Beyond SLG}
\label{sec:orgdcce25c}

Prologos has SLG-style tabling (\texttt{tabling.rkt}). Extensions to explore:
\subsubsection{Subsumptive Tabling}
\label{sec:org4d63287}

Standard variant tabling memoizes exact call patterns. Subsumptive tabling
recognizes that a more general call subsumes a more specific one:

\begin{verbatim}
;; With subsumptive tabling:
;; solve [ancestor :alice ?x] produces answers for all descendants
;; solve [ancestor :alice :carol] can reuse the table entry
;; (because ancestor(:alice, ?x) subsumes ancestor(:alice, :carol))
\end{verbatim}
\subsubsection{Mode-Directed Tabling}
\label{sec:org0a7f71c}

Combine modes with tabling to control which answers are kept:

\begin{verbatim}
;; Keep only the minimum-cost path (mode-directed)
:- table shortest_path(+, +, min)
defr shortest-path
  | [+from : Node] [+to : Node] [-cost : Nat :aggregate min]
\end{verbatim}
\subsubsection{Magic Sets}
\label{sec:orgb5ad14e}

Transform bottom-up evaluation to be goal-directed (top-down efficiency
with bottom-up termination guarantees). This is particularly relevant for
Datalog-style reasoning within Prologos.
\subsection{Implementation Priority Matrix}
\label{sec:orgd6a5b61}

\begin{center}
\begin{tabular}{lllll}
Extension & Infrastructure Reuse & New Lattices & New Syntax & Complexity\\
\hline
Linear LP & QTT multiplicities & Counting & Moderate & Medium\\
Temporal LP & Bilattice (co/ind) & Trace prefix & Moderate & Medium\\
Epistemic LP & ATMS worldviews & None & Light & Low\\
Abductive LP & ATMS diagnoses & None & Light & Low\\
ASP integration & ATMS (amb, nogood) & None & Moderate & Medium\\
CHR & Propagator network & None & Heavy & High\\
Subsumptive tabling & tabling.rkt & None & None & Medium\\
Mode-directed tab. & tabling.rkt + modes & Aggregate & Light & Medium\\
Magic sets & tabling.rkt & None & None & High\\
\end{tabular}
\end{center}
\section{The Frontier: Modern LP Redesign (2020--2026)}
\label{sec:orga65b27f}

\subsection{Survey of Recent Systems}
\label{sec:orgc185050}

\subsubsection{Flix (Aarhus/Waterloo, 2016--2025)}
\label{sec:org276f3b8}

Flix extends Datalog with arbitrary lattices and monotone functions, integrated
as a full programming language. Key innovations:

\begin{itemize}
\item \emph{Lattice semantics}: Datalog relations can map to lattice values, not just sets
of tuples. Computation proceeds to the least fixpoint over the lattice.
\item \emph{First-class Datalog programs}: Functions can take and return Datalog programs
as values (rho abstraction), enabling modular analysis composition.
\item \emph{Type classes}: Flix uses type classes for lattice operations, paralleling
Prologos's trait system.
\item \emph{Stratified negation + lattice fixpoints}: Combines standard Datalog negation
with lattice-valued computation.
\end{itemize}

Relevance to Prologos: Flix demonstrates that lattice-parameterized fixpoint
computation (exactly what Prologos's propagator network does) is a viable
foundation for a full language. Prologos goes further by adding backtracking
search, ATMS, and well-founded semantics.
\subsubsection{Formulog (Harvard, 2020--2024)}
\label{sec:orgc25c6a7}

Formulog = Datalog + SMT + ML-like functions. Key innovations:

\begin{itemize}
\item \emph{SMT integration}: Datalog rules can contain SMT constraints; the solver
dispatches to Z3/CVC5 for satisfiability checking.
\item \emph{Eager evaluation}: A 2024 paper (Bembenek et al., OOPSLA 2024) shows that
DFS-style eager evaluation achieves 5--8x speedups over Soufflé-style
semi-naive evaluation for SMT-heavy workloads.
\item \emph{Algebraic data types}: First-class ADTs in Datalog, enabling analysis of
languages with recursive data.
\end{itemize}

Relevance to Prologos: Formulog's SMT integration could be modeled as an
external constraint domain (like CLP(SMT)). The eager evaluation strategy
is interesting for Prologos's solver, which already uses DFS.
\subsubsection{Ascent (2022--2024)}
\label{sec:org98d8425}

Ascent embeds Datalog in Rust via macros, with:
\begin{itemize}
\item Semi-naive evaluation compiled to Rust
\item Lattice support (similar to Flix)
\item Incremental computation
\end{itemize}
\subsubsection{Crepe (2021--2024)}
\label{sec:org984b4f0}

Crepe provides Datalog as a Rust procedural macro:
\begin{itemize}
\item Semi-naive evaluation with stratified negation
\item Seamless Rust interop
\item Compile-time Datalog analysis
\end{itemize}
\subsubsection{Flan (Purdue, POPL 2024)}
\label{sec:org563e20a}

Flan embeds Datalog in Scala via multi-stage programming:
\begin{itemize}
\item Combines Flix's flexibility with Soufflé's performance
\item Arbitrary aggregates, user-defined functions, lattices
\item Seamless host language integration
\end{itemize}
\subsubsection{Soufflé (2016--2025)}
\label{sec:org6072e13}

High-performance Datalog compiler:
\begin{itemize}
\item Compiles to parallel C++
\item Used in production for program analysis (Java points-to, security)
\item Subsumptive tabling, aggregates, ADTs
\end{itemize}
\subsection{What These Systems Teach Us}
\label{sec:org6df367c}

\begin{center}
\begin{tabular}{lllllll}
Capability & Classic Prolog & Flix & Formulog & Soufflé & Prologos (current) & Prologos (proposed)\\
\hline
Backtracking search & Yes & No & No & No & Yes & Yes\\
Lattice values & No & Yes & No & Partial & Yes (propagator) & Yes\\
SMT integration & No & No & Yes & No & No & Via CLP(SMT)\\
Dependent types & No & No & No & No & Yes & Yes\\
Session types & No & No & No & No & Yes & Yes\\
Linear types (QTT) & No & No & No & No & Yes & Yes\\
Well-founded semantics & XSB & No & No & No & Yes (bilattice) & Yes\\
ATMS/TMS & No & No & No & No & Yes & Yes\\
Probabilistic & ProbLog ext. & No & No & No & No & Yes (semiring)\\
Coinductive & CoLP ext. & No & No & No & No & Yes (bilattice)\\
Modes & Mercury ext. & No & No & No & No & Yes\\
ILP & Metagol ext. & No & No & No & No & Yes (ATMS)\\
Incremental & No & No & No & Partial & No & Via propagators\\
\end{tabular}
\end{center}

The table reveals Prologos's unique position: it is the only system that
combines backtracking search with lattice-based propagators, dependent types,
and ATMS-based assumption management. The proposed extensions would make it
the most expressive logic programming system in existence.
\subsection{The Propagator Network as Universal Substrate}
\label{sec:orgcdaf224}

The central thesis of this document: \emph{every extension of logic programming
maps to a specific configuration of the propagator network}.

\begin{center}
\begin{tabular}{ll}
LP Extension & Propagator Configuration\\
\hline
Standard LP & Boolean cells, unification propagators, DFS scheduler\\
CLP(FD) & Domain cells, arc-consistency propagators\\
Well-founded LP & Bilattice cells, alternating fixpoint scheduler\\
Probabilistic LP & Probability cells (semiring), knowledge compilation via ATMS\\
Coinductive LP & Descending cells (gfp), coinductive hypothesis propagators\\
ASP & ATMS assumptions as choices, nogoods as integrity constraints\\
Abductive LP & ATMS assumptions as abducibles, minimal diagnosis\\
Fuzzy LP & \texttt{[0,1]} cells, fuzzy connective propagators\\
Temporal LP & Trace-prefix cells, unfolding propagators\\
Linear LP & Counting cells, resource-consumption propagators\\
Abstract interp. & Galois-connected cross-domain propagators\\
Datalog/Flix-style & Lattice cells, semi-naive evaluation as propagation strategy\\
\end{tabular}
\end{center}

This is not a coincidence. The Radul-Sussman propagator model is a
\emph{generalization} of constraint propagation, which itself generalizes all of
the above. Prologos's architecture makes this explicit and exploitable.
\section{Composition Map}
\label{sec:orgd8adb2c}

How do these extensions compose with each other and with existing Prologos features?
\subsection{Pairwise Composition Matrix}
\label{sec:org03fda05}

\begin{center}
\begin{tabular}{llllllll}
 & Modes & Domains & Coinduction & ILP & Probabilistic & Category & Linear\\
\hline
\textbf{Modes} & --- & Strong & Moderate & Weak & Moderate & Strong & Strong\\
\textbf{Domains} &  & --- & Weak & Weak & Strong & Strong & Moderate\\
\textbf{Coinduction} &  &  & --- & Weak & Moderate & Strong & Moderate\\
\textbf{ILP} &  &  &  & --- & Strong & Weak & Weak\\
\textbf{Probabilistic} &  &  &  &  & --- & Strong & Moderate\\
\textbf{Category} &  &  &  &  &  & --- & Strong\\
\textbf{Linear} &  &  &  &  &  &  & ---\\
\end{tabular}
\end{center}

Legend: Strong = natural synergy, Moderate = composable with some work, Weak = mostly independent.
\subsection{Key Composition Synergies}
\label{sec:org3505d37}

\subsubsection{Modes + Domains (Strong)}
\label{sec:org2ab4411}

Mode analysis determines propagation direction; domain constraints determine
what propagates. Together: mode \texttt{(+)} on a domain variable means the domain
is fully known at entry; mode \texttt{(-)} means the domain must be computed.
This enables the constraint solver to select arc-consistency (for \texttt{?} mode)
vs. direct evaluation (for \texttt{+} mode).
\subsubsection{Probabilistic + Domains (Strong)}
\label{sec:orga1a9090}

CLP(FD) with probabilities = probabilistic constraint satisfaction. The
probability cell tracks the probability of each domain value, enabling
marginal inference over constrained variables. DC-ProbLog already combines
distributional clauses with constraint domains.
\subsubsection{Probabilistic + ILP (Strong)}
\label{sec:org4de463e}

Learning probabilistic logic programs. ProbLog + ILP = structure learning
for probabilistic models. The ATMS tracks assumption probabilities, and
the gradient semiring enables differentiable structure learning.
\subsubsection{Coinduction + Category Theory (Strong)}
\label{sec:org4ea9560}

Coinductive types are terminal coalgebras; coinductive LP computes greatest
fixpoints. The categorical framework provides the theory (final coalgebra
semantics), and coinductive LP provides the computation. Together: coinductive
reasoning about infinite structures with categorical guarantees.
\subsubsection{Modes + Linear Types (Strong)}
\label{sec:org8308b13}

Mode analysis tells you which arguments are consumed (\texttt{+} in linear context)
and which are produced (\texttt{-}). This determines the resource flow. QTT
multiplicities (\texttt{m0}, \texttt{m1}, \texttt{mw}) already track this at the type level;
modes extend it to the relational level.
\subsubsection{Category Theory + Probabilistic (Strong)}
\label{sec:org7fe1b8a}

The Giry monad provides the categorical foundation for probabilistic
programming. Semiring-based LP is a functor from the category of semirings
to the category of LP semantics. The propagator network is the \emph{universal
construction} that internalizes this functor.
\subsection{Composition with Existing Prologos Features}
\label{sec:orgb401b32}

\subsubsection{With Dependent Types}
\label{sec:org52548e7}

Every extension must respect the type system. Modes become \emph{type-level mode
annotations}. Domain constraints become \emph{refinement types}. Probabilities
become \emph{graded modalities}. The fibered category framework ensures coherence.
\subsubsection{With Session Types}
\label{sec:orgda63157}

Session types are coinductive protocols. Probabilistic session types (Das \&
Pfenning, 2023) track expected resource usage. Linear session types ensure
protocol adherence. Modes on session types determine whether a channel is
for sending (\texttt{+}) or receiving (\texttt{-}).
\subsubsection{With the Propagator Network}
\label{sec:org9c64dad}

All extensions are \emph{propagator network configurations}. The network is the
universal substrate. New extensions add:
\begin{itemize}
\item New cell types (domain cells, probability cells, trace cells, counting cells)
\item New propagator types (arc-consistency, gradient, coinductive hypothesis)
\item New scheduler strategies (semi-naive, alternating fixpoint, belief propagation)
\item New lattice descriptors (semirings, interval lattices, bilattices)
\end{itemize}
\subsubsection{With the ATMS}
\label{sec:org692254c}

The ATMS is the universal \emph{assumption management} layer:
\begin{itemize}
\item Probabilistic LP: assumptions with probabilities
\item Abductive LP: assumptions as abducibles
\item ASP: assumptions as choices, nogoods as constraints
\item ILP: assumptions as candidate clauses
\item Epistemic LP: assumptions as agent observations
\end{itemize}
\section{Research References}
\label{sec:org82d6bf7}

\subsection{Mode Analysis and Optimization}
\label{sec:org44566d9}
\begin{itemize}
\item Somogyi, Henderson, Conway. "The Execution Algorithm of Mercury." JLP, 1996.
\item Overton, Somogyi, Stuckey. "Constraint-Based Mode Analysis of Mercury." PPDP, 2002.
\item Hermenegildo et al. "An Overview of the Ciao Multiparadigm Language and Program Development Environment." ICLP, 2012.
\item Stade et al. "A Practical Approach for Testing Static Analysis Truths." 2025.
\end{itemize}
\subsection{Coinductive Logic Programming}
\label{sec:org55ec01a}
\begin{itemize}
\item Simon, Mallya, Bansal, Gupta. "Coinductive Logic Programming." ICLP, 2006.
\item Gupta, Bansal, Simon. "Coinductive Logic Programming and Its Applications." ICLP, 2007.
\item Ancona, Barbieri, Zucca. "Flexible Coinductive Logic Programming." TPLP, 2020.
\item Ancona, Barbieri, Zucca. "Checking Equivalence of Corecursive Streams." 2024.
\end{itemize}
\subsection{Probabilistic Logic Programming}
\label{sec:org01f4a4e}
\begin{itemize}
\item De Raedt, Kimmig, Toivonen. "ProbLog: A Probabilistic Prolog." UAI, 2007.
\item Manhaeve, Dumancic, Kimmig, Demeester, De Raedt. "DeepProbLog: Neural Probabilistic Logic Programming." NeurIPS, 2018.
\item Derkinderen, Manhaeve et al. "The DeepLog Neurosymbolic Machine." 2025.
\item Zuidberghe et al. "Declarative Probabilistic Logic Programming in Discrete-Continuous Domains." AIJ, 2024.
\item "Semirings for Probabilistic and Neuro-Symbolic Logic Programming." 2024.
\end{itemize}
\subsection{Inductive Logic Programming}
\label{sec:orgf1f89b5}
\begin{itemize}
\item Muggleton, Lin. "Meta-interpretive Learning of Higher-Order Dyadic Datalog." MLJ, 2014.
\item Cropper, Morel. "Learning Programs by Learning from Failures" (Popper). MLJ, 2020.
\item Cropper, Dumancic, Evans, Muggleton. "Inductive Logic Programming at 30." MLJ, 2021.
\item Morel, Cropper. "Typed Meta-interpretive Learning of Logic Programs." ILP, 2019.
\item "Meta-Interpretive Learning with Reuse." Mathematics, 2024.
\item "GLIDR: Graph-Like Inductive Logic Programming with Differentiable Reasoning." 2025.
\item Rocha et al. "Program Synthesis Using ILP for the ARC." 2025.
\item "Bridging Logic Programming and Deep Learning for Explainability through ILASP." 2025.
\end{itemize}
\subsection{Category Theory and Logic}
\label{sec:orgb54bec3}
\begin{itemize}
\item Jacobs. "Categorical Logic and Type Theory." Elsevier, 1999.
\item Goguen, Burstall. "Institutions: Abstract Model Theory for Specification and Programming." JACM, 1992.
\item Hinze. "Kan Extensions for Program Optimisation." MPC, 2012.
\item "Category-Theoretical and Topos-Theoretical Frameworks in Machine Learning: A Survey." Axioms, 2025.
\item Shiebler. "Kan Extensions in Data Science and Machine Learning." 2022.
\end{itemize}
\subsection{Well-Founded Semantics and Bilattices}
\label{sec:org3809080}
\begin{itemize}
\item Fitting. "Bilattices and the Semantics of Logic Programming." FACS, 1991.
\item Van Gelder, Ross, Schlipf. "The Well-Founded Semantics for General Logic Programs." JACM, 1991.
\item Denecker, Marek, Truszczynski. "Approximation Fixpoint Theory." 2004.
\item Charalambidis et al. "The Stable Model Semantics for Higher-Order Logic Programming." 2024.
\end{itemize}
\subsection{Modern Datalog Systems}
\label{sec:orgcf0b7b3}
\begin{itemize}
\item Madsen, Yee, Lhotak. "From Datalog to Flix." PLDI, 2016.
\item Bembenek, Greenberg, Chong. "Formulog: Datalog for SMT-Based Static Analysis." OOPSLA, 2020.
\item Bembenek, Greenberg, Chong. "Making Formulog Fast." OOPSLA, 2024.
\item "Flix: A Design for Language-Integrated Datalog." PACMPL, 2025.
\item Abeysinghe et al. "Flan: An Expressive and Efficient Datalog Compiler." POPL, 2024.
\end{itemize}
\subsection{Linear Logic and Session Types}
\label{sec:orgdf65c8a}
\begin{itemize}
\item Caires, Pfenning. "Session Types as Intuitionistic Linear Propositions." CONCUR, 2010.
\item "Comparing Session Type Systems Derived from Linear Logic." 2024.
\item "The Linear Session Abstract Machine." 2024.
\item Das, Pfenning. "Probabilistic Resource-Aware Session Types." POPL, 2023.
\item "Deadlock-Free Separation Logic: Linearity Yields Progress." POPL, 2024.
\item Derakhshan, Pfenning. "Rast: Resource-Aware Session Types with Arithmetic Refinements." FSCD, 2020.
\end{itemize}
\subsection{ASP and CHR}
\label{sec:org706bbdc}
\begin{itemize}
\item Abdennadher et al. "First-Order ASP Programs as CHR Programs." SAC, 2021.
\item "Quantified Constraint Handling Rules." TOCL, 2025.
\end{itemize}
\subsection{Abductive Logic Programming}
\label{sec:org46f76ee}
\begin{itemize}
\item "An Efficient Propositional System for Abductive Logic Programming." AI Review, 2024.
\end{itemize}
\subsection{Tabling}
\label{sec:orgccec5d4}
\begin{itemize}
\item Swift, Warren. "XSB: Extending Prolog with Tabled Logic Programming." TPLP, 2012.
\item Warren. "Programming in Tabled Prolog." Draft, 2023.
\end{itemize}
\section{Appendix: Lattice Requirements per Extension}
\label{sec:org937b7d0}

\begin{center}
\begin{tabular}{lllll}
Extension & Cell Lattice & Direction & Merge & Needs New \texttt{lattice-desc}?\\
\hline
Standard LP & Boolean & Ascending & Disjunction & No (\texttt{bool-lattice})\\
WFS & Boolean bilattice & Both & Join/Meet & No (\texttt{bilattice-var})\\
CLP(FD) & Reverse powerset & Descending & Intersection & Yes\\
CLP(Bounds) & Interval & Descending & Narrow & Yes\\
Probabilistic & \texttt{[0,1]} semiring & Ascending & Addition & Yes\\
Viterbi & \texttt{[0,1]} max-semiring & Ascending & Max & Yes\\
Fuzzy & \texttt{[0,1]} Heyting & Ascending & Max & Yes\\
Coinductive & Boolean & Descending & Conjunction & No (reuse \texttt{bool-lattice})\\
Linear resource & \texttt{Z} (integers) & Ascending & Addition & Yes\\
Trace prefix & Prefix order & Ascending & Common prefix & Yes\\
Sign abstract & Sign lattice & Ascending & Join & Yes\\
Set interval & Pair of sets & Both & Narrow bounds & Yes\\
\end{tabular}
\end{center}
\section{Appendix: Implementation Priority}
\label{sec:orgfbb5f69}

Based on infrastructure reuse, expressive power gained, and implementation complexity:
\subsection{Tier 1: Low-Hanging Fruit (High Reuse, Moderate Effort)}
\label{sec:org8fe1adf}

\begin{enumerate}
\item \textbf{Abductive LP} --- Almost free: ATMS assumptions as abducibles,
\texttt{atms-minimal-diagnoses} as explanation engine. Needs only surface syntax.

\item \textbf{Epistemic LP} --- ATMS worldviews as agent knowledge states. Surface syntax +
thin dispatch layer over existing \texttt{atms-solve-all}.

\item \textbf{ASP Integration} --- ATMS \texttt{amb} for choice rules, \texttt{nogood} for integrity
constraints. The hardest part is the surface syntax for aggregates.
\end{enumerate}
\subsection{Tier 2: Medium Effort, High Impact}
\label{sec:org7764653}

\begin{enumerate}
\item \textbf{Mode Analysis} --- Significant optimization potential. Requires new analysis
pass but no new propagator infrastructure.

\item \textbf{Coinductive LP} --- Reuses bilattice descending cells. Requires coinductive
hypothesis detection in the DFS solver and coclause syntax.

\item \textbf{Domain-Constrained Variables (CLP(FD))} --- New lattice types but standard
propagator infrastructure. High practical value for constraint satisfaction.
\end{enumerate}
\subsection{Tier 3: High Effort, Transformative Impact}
\label{sec:orgd5cff30}

\begin{enumerate}
\item \textbf{Probabilistic LP} --- New semiring lattices, probability computation, and
potentially knowledge compilation. Very high expressive payoff.

\item \textbf{CHR} --- Full committed-choice rule engine. High implementation cost but
provides a meta-level constraint programming capability.

\item \textbf{ILP/Rule Learning} --- Requires metarule search, hypothesis management,
and potentially differentiable infrastructure.
\end{enumerate}
\subsection{Tier 4: Research-Grade}
\label{sec:org9c5efb8}

\begin{enumerate}
\item \textbf{Category-theoretic topos parameterization} --- Foundational but speculative.
Provides the theoretical framework for all other extensions.

\item \textbf{Institution-based multi-logic dispatch} --- Enables mixing logics in a
single query. Requires significant design work.

\item \textbf{Temporal LP} --- Useful for verification but requires trace infrastructure
and careful interaction with the scheduler.
\end{enumerate}
\end{document}
